\documentclass[10pt,a4paper]{amsart}
\usepackage[margin=19mm]{geometry}
\usepackage{amsmath,amssymb,mathtools}
\usepackage[scr=rsfs,cal=euler]{mathalfa}
\usepackage{xfrac}
\usepackage[unicode,hidelinks,bookmarks=false]{hyperref}
\newcommand{\ie}{i.e.\ }
\newcommand{\cf}{cf.\ }
\newcommand{\resp}{respectively\ }
\newcommand{\Subsection}{\subsection}
\providecommand{\nobreakdash}{\nobreak\hbox{-}}
\setlength{\emergencystretch}{2em}
\multlinegap0pt
\usepackage{bm}
\usepackage{bbm}
\usepackage{dsfont}
\usepackage{xspace}
\usepackage{cancel}
\usepackage{mathtools}
\usepackage{amsthm}
\makeatletter
\@ifundefined{lemma}{\newtheorem{lemma}{Lemma}}{}
\@ifundefined{theorem}{\newtheorem{theorem}{Theorem}}{}
\makeatother
\makeatletter
\newcommand{\bP}{\mathbb{P}}
\newcommand{\cP}{\mathcal{P}}
\expandafter\ifx\csname customproofsketch\endcsname\relax
\newenvironment{customproofsketch}{\noindent\emph{Proof} (sketch).}{\qed}
\fi
\makeatother
\title[Robust Competitive Ratio for Deterministic Monopoly Pricing]{Robust Competitive Ratio for Deterministic Monopoly Pricing}
\author{Tim S.G. van Eck, Pieter Kleer, Johan S.H. van Leeuwaarden}
\date{Source: arXiv:2509.06619v1 (CC BY 4.0)}
\begin{document}
\maketitle

\noindent\emph{Source and licence.} Robust Competitive Ratio for Deterministic Monopoly Pricing. Version arXiv:2509.06619v1 (CC BY 4.0), distributed under the Creative Commons Attribution 4.0 International licence, as recorded in the arXiv licence metadata for this exact version and shown on the abstract page https://arxiv.org/abs/2509.06619v1. Full rights notice: licenses/arxiv-2509.06619-CC-BY-4.0.txt.\par\smallskip

\noindent\textit{Editorial scope.} Counted unit: the Lemma whose source label is \texttt{None} in the article (source0.tex lines 1046--1048, source label \texttt{None}), delivered together with the article's own complete original proof of it (source0.tex lines 1049--1071, one \texttt{proof} environment of 23 lines). No mathematics is written, repaired or re-derived: statement and proof are the article's own text line for line. Required context retained uncounted: th:optimal\_price\_variance (lines 634--649); lemma:price1 (lines 1005--1007). Every definition, display and label of those lines is retained unchanged. Omitted from this package and not redistributed: 1--633, 650--1004, 1008--1045, 1072--1114. Nothing inside a retained excerpt is omitted or altered: every retained body is delivered exactly as the article's lines are written, so no omission receipt of any kind accompanies this package. Citations: the retained excerpts carry no citation command at all, so no bibliography entry of the article is transcribed and no reference is redistributed. Reference adaptation: none was needed and none was made; every reference inside the delivered unit resolves inside it, so no unresolved reference or citation is produced. Printed slips kept literal: no printed word, symbol, hypothesis or number of the counted statement or of its proof was altered. Layout and class: the article's own class is \texttt{article} with packages including amsmath, graphicx, subcaption, caption, geometry, booktabs, lscape, multirow, fontenc, babel, a4wide, amssymb, xcolor, inputenc; the wrapper is the lane's pinned \texttt{amsart} editorial wrapper at 10pt on A4, which restates the theorem-like environments the retained text needs and adds the article's own macro definitions that the retained text needs (\texttt{\textbackslash bP}, \texttt{\textbackslash cP}), with extra wrapper packages (bm, bbm, dsfont, xspace, cancel, mathtools). The delivered numbering is the unit's own. Assets: no retained span contains a figure environment, a graphics-inclusion command, a PDF-page inclusion, a table or a TikZ picture; no image file is copied into this package, the package declares no asset dependency and no image file is redistributed with it. Licence: version arXiv:2509.06619v1 of this article is distributed under the Creative Commons Attribution 4.0 International licence, as recorded in the arXiv licence metadata for this exact version and shown on the abstract page https://arxiv.org/abs/2509.06619v1. A full rights notice is delivered at licenses/arxiv-2509.06619-CC-BY-4.0.txt. Characters: every retained excerpt is pure ASCII, so the delivered engine, XeLaTeX with its default Latin Modern text font, reports no missing character.
\begin{theorem}\label{th:optimal_price_variance}
    The optimal robust price $p_{(\mu,\sigma,\beta)}^*=\arg\sup_{p>0}\inf_{\bP\in\cP(\mu,\sigma,\beta)}\textup{CR}(p,\bP)\nonumber$ is
    \begin{align}
        p_{(\mu,\sigma,\beta)}^* = \begin{cases}
            p_l^*, \quad & \sigma \leq \sigma^{*}, \\
            p_h^* = \max\{p_{h1}^*,p_{h2}^*\}, \quad &    \sigma \geq \sigma^{*},
        \end{cases}
    \end{align}
with
\begin{align}
    p_{l}^*&=\mu - \sigma \cdot \left(\left(\frac{\mu}{2\sigma}+\sqrt{\frac{8}{27}+\frac{\mu}{2\sigma}}\right)^{\frac{1}{3}}+\left(\frac{\mu}{2\sigma}-\sqrt{\frac{8}{27}+\frac{\mu}{2\sigma}}\right)^{\frac{1}{3}}\right),\\
    p_{h1}^* &= \frac{1}{2}\left(\beta+\tau_2-\sqrt{(3\beta-\tau_2)^2-4\beta^2}\right),\ \ p_{h2}^* = \frac{\tau_2}{2}.
\end{align}
Furthermore, $\sigma^*$ is implicitly defined as the solution to
$R(p_l^*,\sigma^*) = R(p_h^*,\sigma^*)$. 
\end{theorem}

\begin{lemma}\label{lemma:price1}
    The optimal robust price $p^*_{(\mu,\sigma,\beta)} \in \{p^*_l,p^*_h\}$.
\end{lemma}

\begin{lemma}
    Consider $\sigma^*$ from \textup{Theorem \ref{th:optimal_price_variance}}. When $\sigma \leq \sigma^*$, then $p^*_{(\mu,\sigma,\beta)} = p^*_l$ and when $\sigma \geq \sigma^*$, then $p^*_{(\mu,\sigma,\beta)} = p^*_h$.
\end{lemma}

\begin{proof}
    As motivated in the proof of Lemma \ref{lemma:price1}, we will assume without loss of generality that $\mu=1$. Furthermore, we recall that $$f_{1a}(p) = \frac{(1-p)^2}{(1-p)^2+\sigma^2},\ \ f_{1b}(p) = \frac{p(1-p)}{(1-p)+\sigma^2},\ \ f_{2a}(p) = \frac{p(1+\sigma^2-p)}{(\beta-p)(\beta+p-1-\sigma^2)}, \ \ f_{2b}(p) = \frac{p}{\beta}.$$

    Now consider the case that $\sigma^2 \xrightarrow{} \mu(\beta-\mu)$, which is the maximal value of $\sigma^2$. Then $p^*_{(\mu,\sigma,\beta)} = p^*_h$, as $\tau_1 \xrightarrow{} 0$ implies $p^*_l \xrightarrow{} 0$. Next, consider the case that $\sigma^2 = 0$. Then $p^*_{(\mu,\sigma,\beta)} = p^*_l$, as $p^*_l = 1$ with $f_{1a}(p^*_l) = f_{1b}(p^*_l) = 1$, which is the highest value for the competitive ratio, while $p^*_h$ performs strictly worse.
    
    Since $f_{1a}(p)$ and $f_{1b}(p)$ are continuously differentiable in $p^*_l$, we can use the envelope theorem to show that
    \begin{align}
        \frac{\partial f_{1a}(p^*_l)}{\partial \sigma^2} = \frac{\partial f_{1a}(p)}{\partial\sigma^2}\biggr\rvert_{p = p^*_l} = -\frac{\left(1 - p^*_l\right)^{2}}{\left(\sigma^2 + \left(1 - p^*_l\right)^{2}\right)^{2}} \leq 0,
    \end{align}
    and
    \begin{align}
        \frac{\partial f_{1b}(p^*_l)}{\partial \sigma^2} = \frac{\partial f_{1b}(p)}{\partial\sigma^2}\biggr\rvert_{p = p^*_l} = \frac{\left(p^*_l - 1\right) p^*_l}{\left(\sigma^2 - p^*_l + 1\right)^{2}} \leq 0,
    \end{align}
    as $p^*_l \leq 1$. Hence, $\min\{f_{1a}(p^*_l),f_{1b}(p^*_l)\}$ is decreasing in $\sigma^2$. Furthermore, since $f_{2a}(p)$ and $f_{2b}(p)$ are continuously differentiable in $p^*_h$, we can use the envelope theorem to show that
    \begin{align}
        \frac{\partial f_{2a}(p^*_h)}{\partial \sigma^2} = \frac{\partial f_{2a}(p)}{\partial\sigma^2}\biggr\rvert_{p = p^*_h} = \frac{\beta p^*_h}{(\beta-p^*_h)(\sigma^2-p^*_h-\beta+1)^2} \geq 0,
    \end{align}
    and
    \begin{align}
        \frac{\partial f_{2b}(p^*_h)}{\partial \sigma^2} = \frac{\partial f_{2b}(p)}{\partial\sigma^2}\biggr\rvert_{p = p^*_h} = 0 \geq 0.
    \end{align}
    Hence, $\min\{f_{2a}(p^*_h),f_{2b}(p^*_h)\}$ is increasing in $\sigma^2$. This finishes the proof.
\end{proof}

\end{document}
